% =====================================================================
%  c01a.tex —— 原文 pp.1–2（Lead）
%  内容：摘要段、§1 向量丛与 K(X)、定义 1–2、命题 1
% =====================================================================
% ===== 原文 pp.1–2 =====
% [原文 p.1]
% [Tutor Note: 摘要中的 “$\Omega$-prespectrum” 指以环路空间函子 $\Omega$ 衔接的预谱（$E_n\to\Omega E_{n+1}$），
% 第 3 节定义 6 会给出严格定义；此处先留作动机。原文 “the only spheres that are $H$-spaces are $S^k$ for $k=0,1,3,7$”
% 中的 $H$-空间 = 带连续乘法（不必有单位、结合、逆）的空间，见第 4 节定理 6。]

拓扑 $K$ 理论的基本想法是：空间可以凭它所承载的向量丛来区分。下面我们先介绍基本的概念与定义
（向量丛、分类空间），再陈述一条主要定理（Bott 周期性），最后借助 $\Omega$-预谱（$\Omega$-prespectrum）
的可表示性，把 $K$ 理论放进广义上同调的框架。之后我们指出一个应用（仅有的 $H$-空间球面是
$S^k$，$k=0,1,3,7$），以及 $K$ 理论可以走向的其他方向（代数 $K$ 理论与 $C^*$-代数的 $K$ 理论）。

\section{向量丛与 $K(X)$}

\begin{defn}{1}
  向量丛（vector bundle，$k=\R$ 或 $\C$ 上维数为 $n$ 的）由底空间（base space）$B$、全空间
  （total space）$E$ 与映射 $E\xrightarrow{p}B$ 构成，且满足以下性质。
  \begin{itemize}[leftmargin=1.6em,itemsep=2pt]
    \item $E$ 局部是乘积：存在 $B$ 的开覆盖 $\bigcup_i U_i$ 与同胚
      \[ \phi : p^{-1}(U_i) \xrightarrow{\ \cong\ } U_i\times k^n , \]
      使得 $(p\circ\phi_i^{-1})(u,v)=u$。
    \item 转移映射是线性的：同胚
      \[ \phi_j\circ\phi_i^{-1}\big|_{(U_i\cap U_j)\times k^n} :
         (U_i\cap U_j)\times k^n \longrightarrow (U_i\cap U_j)\times k^n \]
      在每一纤维 $\{b\}\times k^n$ 上都是线性同构。
  \end{itemize}
\end{defn}
% [Tutor Note: 定义 1 第一条里原文写作 “there is a an open cover”，冠词重复（原文笔误，照译未改）。
% 另：第一条的展示式印作 $\phi$（不带下标），而紧随其后的公式与第二条分别用 $\phi_i^{-1}$、$\phi_j\circ\phi_i^{-1}$；
% 按上下文展示式应写作 $\phi_i$，原文此处漏了下标。正文照原样排（展示式仍作 $\phi$），此处注明。]

人们也可以用粘合函数（clutching functions）如下地定义 $X$ 上的向量丛。设 $X=\bigcup_i U_i$ 为开覆盖，
$E_i=U_i\times\C^n$ 为覆盖的每个成员上的平凡丛。假设有 $g_{ij}:U_i\cap U_j\to\GL_n(\C)$ 在
$U_i\cap U_j\cap U_k$ 上满足上链条件（cocycle condition）$g_{jk}\circ g_{ij}=g_{ik}$。则
$\bigsqcup_i E_i/\sim$ 是 $X$ 上的向量丛（投影到第一个坐标），其中
\[ (x,v)\sim(x,g_{ij}v),\qquad x\in U_i\cap U_j . \]
% [Tutor Note: 这个商空间要成为丛，须默认 $g_{ii}=1$、$g_{ij}^{-1}=g_{ji}$；上链条件 $g_{jk}g_{ij}=g_{ik}$
% 正是保证等价关系可传递、且转移函数在 $U_i\cap U_j$ 上相容。原文没有明说，补于此。]

一些例子：
\begin{itemize}[leftmargin=1.6em,itemsep=2pt]
  \item 平凡丛 $E=B\times k^n$；
  \item 光滑流形的切丛与浸入流形的法丛；
  \item $P^1(\C)$ 上的典范丛（过 $\C^2$ 原点的直线），
    \[ H=\{(L,v) : L\in P^1(\C),\ v\in L\}. \]
% [原文 p.2]
  \item Möbius 丛 $[0,1]\times\R/\sim$，其中把 $(0,x)$ 与 $(1,-x)$ 等同（这是
    $S^1\cong P^1(\R)$ 上的典范丛）。
\end{itemize}

$P^1(\C)$ 上的典范丛以后会很重要。

给定底空间上的 $k$-向量丛，可以定义加法运算与乘法运算。设同一底空间上有两个 $k$-向量丛
$E\xrightarrow{p}B$ 与 $E'\xrightarrow{p'}B$（维数分别为 $n$ 与 $m$），我们有
\[
\left\{
\begin{array}{l}
E\oplus E'\xrightarrow{p\oplus p'}B\\[2pt]
(p\oplus p')^{-1}(b)=p^{-1}(b)\oplus p'^{-1}(b)
\end{array}
\right.
\qquad
\left\{
\begin{array}{l}
E\otimes E'\xrightarrow{p\otimes p'}B\\[2pt]
(p\otimes p')^{-1}(b)=p^{-1}(b)\otimes p'^{-1}(b)
\end{array}
\right.
\]
即向量丛的逐纤维和与逐纤维积（维数分别为 $n+m$ 与 $nm$）。

若 $B=\bigcup_i U_i$ 使得 $E,E'$ 在 $U_i$ 上平凡，且 $\{g_{ij}\}$、$\{g'_{ij}\}$ 分别是它们的粘合函数，
则 $E\oplus E'$ 与 $E\otimes E'$ 的粘合函数为
\[ \{g_{ij}\oplus g'_{ij}\in\GL_{n+m}(k)\},\qquad \{g_{ij}\otimes g'_{ij}\in\GL_{nm}(k)\}. \]

这些运算使 $B$ 上的 $k$-向量丛的同构类 $\Vect_k(B)$ 具有交换半环的结构。我们可以形式地扩充到
虚拟丛（virtual bundles），这类似于从 $\N$ 构造 $\Z$。

\begin{prop}{1}
  给定半环 $A$，存在环 $R$ 与同态 $A\to R$，使得对任意从 $A$ 到环 $R'$ 的同态 $A\to R'$，
  都存在唯一的同态 $R\to R'$ 使
  \begin{center}
  \begin{tikzcd}
    A \arrow[r] \arrow[dr] & R' \\
    & R \arrow[u]
  \end{tikzcd}
  \end{center}
  交换。此外，这一构造关于 $A$ 是函子的。[由下述构造从阿贝尔半群 $A$ 得到的阿贝尔群 $G(A)$，
  称为 $A$ 的 \textbf{Grothendieck 群}。]
\end{prop}

\begin{proof}
  令 $R$ 为集合 $A\times A$ 模去如下等价关系：$(a,b)\sim(a',b')$ 当且仅当存在 $c\in A$ 使
  \[ a+b'+c=a'+b+c . \]
  加法与乘法定义为
  \[ (a,b)+(c,d)=(a+b,c+d),\qquad (a,b)\cdot(c,d)=(ac+bd,ad+bc). \]
  另一种构造是取 $A$ 上的自由阿贝尔群，再模去适当的关系
  \[ \{(a\oplus a')-(a+a') : a,a'\in A\}. \]
\end{proof}
% [Tutor Note: p.2 命题 1 证明中的加法公式原文印作 $(a,b)+(c,d)=(a+b,c+d)$，应为 $(a+c,b+d)$
% （判据：$(a,b)$ 表示形式差 $a-b$，加法必须逐分量相加，才能与半群运算相容；且它与下面给出的
% $(a,b)\cdot(c,d)=(ac+bd,ad+bc)$ 一致，后者恰是 $(a-b)(c-d)$ 的分解）。正文照原文排，此处指出。
% 顺带记录原文四处拼写：homorphism → homomorphism、constructin → construction、constuction → construction、isomophism → isomorphism。]

这个虚拟丛构成的环就是我们的研究对象。

\begin{defn}{2}
  复数虚拟丛在 $X$ 上的同构类构成的环称为 $X$ 的 $K$ 理论（$K$-theory），记作 $K(X)$。
  对应 $X\mapsto K(X)$ 是从拓扑空间范畴到环范畴的反变函子。[另有 $X$ 的实 $K$ 理论，
  记作 $KO(X)$。]
\end{defn}
